\chapter{Background Estimation}%
\label{sec:bkg-estimation}

The analysis consists of several backgrounds. The \Wjets, other vector boson
(\Zjets and diboson events, named \emph{Other V/VV} in the following plots and in the tables) and other top (\(s\)-channel and \(t\)-channel
single top, \ttV, \ttH, named \emph{Other top} in the following plots and in the tables) backgrounds are estimated using Monte Carlo
simulation. Additionally the normalisation of the \Wjets background is
constrained in the fit described in \cref{sec:WjetsCR}.

Only prompt background leptons originating from decays of \Zboson, \Wboson, and
\Hboson bosons, or from prompt leptonic \( \tau \) decays are of interest for the analysis.
Fake lepton contributions come from processes such
as in-flight decays of mesons, jets reconstructed as leptons, and electron
photon conversions, this background is named \emph{Misidentified leptons} in the following plots and in the tables.

Prompt leptons are estimated with Monte Carlo and identified with truth-reconstruction
object matching to remove the overlap with the fake-factor method described
in \cref{sec:bkg-fakefactor}. Events containing at least one fake lepton
are estimated with the data-driven fake-factor method.

\section{\texorpdfstring{\Wjets}{W+jets} control region}%
\label{sec:WjetsCR}
The dominant background in the analysis arises from \Wjets background processes.
In order to effectively determine the background estimate and constrain
the normalisation of the \Wjets background across all measurement regions,
a dedicated control region is defined. This control region is designed to have
similar kinematic properties as the respective SRs but involves an inverted
variable selection. This inversion effectively reduces signal contamination
while increasing the yield and purity of a specific background.

The CR for the \Wjets background contains exactly three jets, which are selected
to ensure orthogonality to the SR.\@ At least one of the jets is \btagged.
Additionally the \Wboson-hadronic signal region is vetoed although the overlap is expected
to be small.
To enhance the fraction of \Wjets events, a cut centred around the \Wboson boson
mass is applied through the variable \mtW. Additionally, to reduce the \ttbar
contribution, a cut on \(\mblmin > \qty{180}{\GeV}\) is set.
The resulting \Wjets selection is shown in \cref{tab:RegionDefinitionsCR}
and the corresponding yields in \cref{tab:CR_Wsys:yields}.

\begin{table}[htbp]
  \begin{center}
    \begin{tabular}{l c}
      \toprule
      Selection              & \Wjets CR            \\
      \midrule
      \( N \)(central jet)   & \( = 3 \)            \\
      \( N_{\bjets} \)       & \( \geq 1 \)         \\
      \( \mtW + \met \)      & \( (60, \infty) \)   \\
      \( \pt(\ell) \)        & \( (30, \infty) \)   \\
      \( \pt(j)\)            & \( (25, \infty) \)   \\
      \( \mtW \)             & \( (40, 100) \)      \\
      \mblmin                & \( (180, \infty) \)  \\
      SR veto                & enabled              \\
      \bottomrule
    \end{tabular}
  \end{center}
  \caption{A summary of the event selection for the \Wjets control region.}%
  \label{tab:RegionDefinitionsCR}
\end{table}

\begin{table}[htbp]
  \begin{center}
    \input{tables/yields_CR_Wsys_l_3j_excl}
  \end{center}
  \caption{Yields of individual signal and background samples in the
    \Wjets control region.
    The total expected yield is using the nominal \tW DR sample.
    Uncertainties are statistical only.}%
  \label{tab:CR_Wsys:yields}
\end{table}


\cref{fig:WJetsCR:mblmin} shows the distribution of the \mblmin variable.
The \Wjets purity is roughly \qty{50}{\%} and there is no significant
sensitivity to the \ttbar and \tW interference observed.
A single-bin distribution is used in the likelihood fit as described
in \cref{sec:LikelihoodFit} and illustrated in \cref{fig:WJetsCR:nleptons}.

\begin{figure}[htbp]
  \centering
  \subfloat[]{
    \includegraphics[page=1, width=0.37\textwidth]{CR_WJets/CR_Wsys_l_3j_excl}%
    \label{fig:WJetsCR:mblmin}
  }
  \subfloat[]{
    \includegraphics[page=5, width=0.37\textwidth]{CR_WJets/CR_Wsys_l_3j_excl}%
    \label{fig:WJetsCR:nleptons}
  }

  \caption{Distributions in the \Wjets CR.\@
    \protect\subref{fig:WJetsCR:mblmin} \mblmin distribution and
    \protect\subref{fig:WJetsCR:nleptons} \(N_\mathrm{leptons} \) distribution.
    The full uncertainty without any fit taken into account is shown for the nominal setup using the \tW DR sample.
  }%
  \label{fig:WJetsCR}
\end{figure}


\subsection{Likelihood Fit}%
\label{sec:LikelihoodFit}

\subsubsection{Background Information}

The fitting procedure of this analysis is implemented using the
\textsc{HistFitter}~\cite{Baak:2014wma} fitting framework.
A background-only likelihood fit is performed for the dedicated Control Region
(\Wjets CR) introduced in \cref{sec:WjetsCR}.
The fit makes use of a probability density function (PDF) that combines
Poisson terms for all events corresponding
to each background process in all relevant regions.

The normalisation of the \Wjets background is regulated by a free-floating
parameter, also known as a normalisation parameter.
Other parameters, the so-called nuisance parameters are modelled as having Gaussian distributions, which capture uncertainties associated with various sources such as statistical fluctuations and systematic effects. The likelihood fit aims to maximize the logarithm of the likelihood
function, taking into account the observed data in the Control Region and
the expected background predictions.
The pre- and post-fit distributions for the control region are shown in \cref{fig:LikelihoodFitCR}. 

The extrapolation from the CR to the measurement regions (SR) is determined by
the so-called transfer factor for the normalised \Wjets background process.
This factor determines the estimation of the yields in the measurement regions
and is calculated as follows: 

\begin{equation}
  N\mathrm{(SR, est.)} = N\mathrm{(CR, fit.)} \times \dfrac{\mathrm{MC}\mathrm{(SR, init.)}}{\mathrm{MC}\mathrm{(CR, init.)}} = \mu \times \mathrm{MC}\mathrm{(SR, init)}
  \label{eq:transferFactor}
\end{equation}

where \(N\mathrm{(SR, est.)} \) represents the estimated background
in the measurement regions (SR) for the normalised \Wjets background process,
\(N\mathrm{(CR, fit.)}\) is the fitted number of  events for \Wjets in the CR,
and \({\mathrm{MC}\mathrm{(SR, init.)}}\) and \({\mathrm{MC}\mathrm{(CR, init.)}}\)
are the initial estimates of the contributions from the \Wjets background
in the SR and CR respectively, as obtained in MC simulation.
\(N\mathrm{(CR, fit.)}\) is often
rewritten as \( \mu \) multiplied by a fixed, nominal background prediction,
where \( \mu \) is the actual normalisation factor obtained in the fit to data.
The application of the TF allows systematic uncertainties in the predicted
background processes to be (partially) cancelled in the extrapolation.


\subsubsection{Fit Setup}
The analysis employs a straightforward 1-bin fit focusing on a single parameter of interest: the W+jets normalisation. All \Wjets systematics (modelling and detector) are normalised to the nominal sample yield. Internally, The HistFitter framework splits each uncertainty nuisance parameter into a normalisation and shape component for \Wjets. However, for this 1-bin control region, the normalisation component is removed, simplifying the fit. All uncertainties (both shape and normalisation) from other processes are also included in the fit as well, which helps to provide a more accurate estimation of the \Wjets background, especially given the large contamination from \ttbar events.

Initially, only the uncertainties directly affecting the shape of the \Wjets background were included in the fit. However, since the impact on the unfolding results was minimal between this setup and the one with all systematics included, the nominal fit was ultimately performed with the full set of systematic uncertainties. For more details, see \cref{app:FitSystematicsUsage}.
% The framework internally splits each uncertainty nuisance parameter into
% normalisation and shape components. 
% All systematics related to the \Wjets sample are treated in a special way where the normalisation component is disregarded and the systematic samples are normalised to the nominal yield in the control
% region. This reduces the complexity of the fit and removes unnecessary
% correlations between the normalisation parameter and individual systematics
% nuisance parameters.  Apart from that, systematics are considered for all samples with normalisation and shape components
% as described in \cref{ch:systematics}. 


\subsubsection{Results}


The obtained normalisation factor for \Wjets background from the fit to data is 1.095 \(\pm \) 0.099.

\cref{fig:LikelihoodFit:pullPlot} shows the pull plots of individual systematic
nuisance parameters. The central values are consistently zero and none of the
parameters is over- or under-constrained.

\begin{figure}[htbp]
  \centering
  \includegraphics[page=1, width=1.0\textwidth]{LikelihoodFit/pull_WbWb_WFit_sysBKG_bulk_sigdefault}%
  \caption{Pull plot for CR-only likelihood fit. Luminosity point actually does not represent a pull but is automatically put on this plot by \textsc{HistFitter}.}%
  \label{fig:LikelihoodFit:pullPlot}
\end{figure}

The correlation matrix \cref{fig:LikelihoodFit:corrMatrix} displays only correlations between the normalisation factor for \Wjets background and fake electrons and some signal (\ttbar and single top) modelling uncertainties (ISR, renormalisation scale).
The shown 6 nuisance parameters have the highest (anti-)correlations, the rest % chktex 36
are negligible and not shown. Due to the use of 1-bin fit by construction all
nuisance parameters that are shared between the \Wjets sample and any other sample
do not impact the normalisation significantly as they are collectively constrained
by data. On the other hand nuisance parameters that are not shared with the \Wjets
sample are directly (anti-)correlated, for example cross-section uncertainty on the signal: % chktex 36
when signal increases the floating \Wjets sample has to decrease to maintain the same
constant yield as data has.

\begin{figure}[htbp]
  \centering
  \includegraphics[page=1, width=1.0\textwidth]{LikelihoodFit/c_corrMatrix_RooExpandedFitResult_afterFit}%
  \caption{Correlation matrix for CR-only likelihood fit. ``Normalisation\_signal'' parameter refers to the scale uncertainties for the single top and \ttbar signals}%
  \label{fig:LikelihoodFit:corrMatrix}
\end{figure}


\begin{figure}[htbp]
  \centering
  \includegraphics[width=0.5\textwidth]{LikelihoodFit/WjetsRanking.png}%
  \caption{Ranking of uncertainties for the CR-only likelihood fit}
\end{figure}

The distributions without and with fit results taken into account for the bulk and the search-like
measurement regions (defined in \cref{ch:regions}) are shown in \cref{fig:LikelihoodFitBulk,fig:LikelihoodFitSearch},
respectively. A small reduction in uncertainties is observed in the tail
of the \mblmin and \amttwo variables in the bulk measurement region where
the fraction on \Wjets background is particularly high.
The uncertainty is still high as the signal modelling uncertainty is unchanged
and only the \Wjets background is affected by the fit.

\begin{figure}[htbp]
  \centering
  \subfloat[]{
    \includegraphics[width=0.37\textwidth]{figures/LikelihoodFit/CR_Wsys_l_3j_excl_nleptons_log_pre}%
    \label{fig:LikelihoodFitCR:BulkPreFit}
  }
  \subfloat[]{
    \includegraphics[width=0.37\textwidth]{figures/LikelihoodFit/CR_Wsys_l_3j_excl_nleptons_log_post}%
    \label{fig:LikelihoodFitCR:BulkPostFit}
  }
  \caption{Pre-fit and post-fit distributions for the lepton multiplicity in the \Wjets control region.}%
  \label{fig:LikelihoodFitCR}
\end{figure}

\begin{figure}[htbp]
  \centering
  \subfloat[]{
    \includegraphics[width=0.37\textwidth]{figures/LikelihoodFit/IR/mblmin_nofit}%
  }
  \subfloat[]{
    \includegraphics[width=0.37\textwidth]{figures/LikelihoodFit/IR/mblmin_fit}%
  }
  \caption{Distributions of \mblmin in the bulk measurement region without fit inputs (left) and with fit inputs (right).}%
  \label{fig:LikelihoodFitBulk}
\end{figure}

\begin{figure}[htbp]
  \centering
  \subfloat[]{
    \includegraphics[width=0.37\textwidth]{figures/LikelihoodFit/SR/mblmin_nofit}%
  }
  \subfloat[]{
    \includegraphics[width=0.37\textwidth]{figures/LikelihoodFit/SR/mblmin_fit}%
  }

  \caption{Distributions of \mblmin in the search-like measurement region without fit inputs (left) and with fit inputs (right).}%
  \label{fig:LikelihoodFitSearch}
\end{figure}

\begin{figure}[htbp]
  \centering
  %\includegraphics[width=0.4\textwidth]{LikelihoodFit/CorrTestSystematic.png}%
  \includegraphics[width=0.4\textwidth]{LikelihoodFit/ratioHistBKGRelDiff.pdf}
  \caption{Test of impact of correlations in systematic uncertainties on the \Wjets background modelling. The plot shows the impact of the shifted mu value on the background modelling systematics.}%
  \label{f:testcorrelation}
\end{figure}

\subsubsection{Propagation of uncertainties to unfolding}
In the unfolding procedure, uncertainties are treated as uncorrelated. From the fit, the individual normalisation parameter of the \Wjets background, its uncertainty, and the shape-only component of other \Wjets uncertainties are propagated to the unfolding.
Given that the normalisation parameter incorporates correlation with the normalisation of the other systematic uncertainties, using it in the unfolding may introduce a slight over- or underestimation of the propagated uncertainty. While the correlations within \Wjets uncertainties are minimal and have a negligible impact—since the normalisation parameter primarily reflects the normalization component of \Wjets systematic uncertainties—this assumption does not hold for modelling uncertainties from other samples.
With this in mind, the correlation matrix \cref{fig:LikelihoodFit:corrMatrix} shows strong (anti-)correlations between the \Wjets normalisation parameter and specific uncertainties, particularly from fake electrons and ISR signal modelling. By propagating the normalisation parameter with these embedded correlations to the unfolding, we introduce a minor double-counting effect, leading to a more conservative (overestimated) uncertainty on the total result. % chktex 36
This effect was roughly validated by re-fitting the mu \Wjets parameter after moving the nominal value of the dominant systematic uncertainty (top shower modelling systematic) to its +1 sigma variation. The resulting mu parameter is then used through the unfolding to assess a new version of the systematic under study. The systematic is not changed by this procedure and the two systematics are identical. The new mu value has an impact on the \Wjets background modelling systematic which is also shown in \cref{f:testcorrelation}. 





\FloatBarrier

%-------------------------------------------------------------------------------

\section{Fake-factor Method}%
\label{sec:bkg-fakefactor}

Fakes for electrons and muons are estimated with the so-called
\textit{fake-factor method}.
The fake factor method is an extrapolation technique for electron or muon fakes
which are taken from a control region rich in this kind of background and with
high statistical power. A weight is calculated for each fake lepton called
the fake factor which is related with the misidentification probability
for a fake lepton to satisfy the selection requirement of a prompt one. Events
containing fake leptons are multiplied with these weights.
Fake factors are measured in a control sample in data orthogonal with respect
to the selections of the main analysis.
These control samples are defined for electrons in \cref{sec:bkg-fakes-electron}
and for muons in \cref{sec:bkg-fakes-muon}.
The misidentification probability, sometimes referred to as \textit{fake rate}
is defined as
\begin{equation}
  f = \dfrac{N_\mathrm{pass}}{N_\mathrm{pass}+N_\mathrm{fail}}
  \label{eq:fake-rate}
\end{equation}
where \( N_\mathrm{pass} \) and \( N_\mathrm{fail} \) are the numbers of fake
leptons satisfying and failing the fake identification requirement. The fake factor
is defined as
\begin{equation}
  F = \frac{f}{1-f} =  \dfrac{N_\mathrm{pass}}{N_\mathrm{fail}}
  \label{eq:fake-factor}
\end{equation}

Once \(F\) is measured as a function of some kinematic variable, it is applied
to fakes taken from data in other analysis regions.
Leptons passing the signal identification
criteria are called \textit{tight} (T) and those failing them \textit{loose} (L).
All criteria are reported in \cref{tab:definitions-electrons,tab:definitions-muons}.
The final amount of extrapolated fake background is defined as
\begin{equation}
  N_\mathrm{fakes} =
  {\left[
  \sum\limits_\mathrm{events}^{n\, \mathrm{loose}}
  {(-1)}^{n-1} \prod\limits_{i}^{n} F_i
  \right]}_\mathrm{data}
  -
  {\left[
  \sum\limits_\mathrm{events}^{n\, \mathrm{loose}}
  {(-1)}^{n-1} \prod\limits_{i}^{n} F_i
  \right]}_\mathrm{prompt~simulation}
\end{equation}
where the inner product runs over the number of loose leptons in the event and
the outer sum over all events with at least one loose lepton present.
Tight leptons do not contribute to the final number of fake leptons.
To avoid double counting, the prompt contribution estimated with MC events needs
to be subtracted from data, as indicated with square brackets in the formula.
The derivation of the formulae are presented in more detail
in Ref.~\cite{EGAM-2019-01}.

In case of single lepton events the formula simplifies into
\begin{equation}
  N^\mathrm{fakes} =
  {\left[ \sum\limits_{L}F \right]}_\mathrm{data} -
  {\left[ \sum\limits_{L}F \right]}_\mathrm{prompt~simulation}
\end{equation}

%-------------------------------------------------------------------------------

% Fake Background Electrons
\section{Electron Fake Factor Measurement}%
\label{sec:bkg-fakes-electron}

Fake electrons arise from three background sources. They are either
mis-reconstructed jets, mostly originating from light quark or gluon
fragmentation, photon conversions or real electrons arising from secondary
decays of light and heavy flavour mesons within jets. Fake electrons are
normally less isolated than prompt electrons and have a lower chance of passing
tighter identification levels. The simulation does not offer an optimal
modelling of fake electrons which ultimately arise as the final product of
a chain prompted by strong interactions.

The electron fake contribution to the measurement region is estimated as described
in \cref{sec:bkg-fakefactor}, which involves a data driven estimate of
the fake factor.

Electron fake factors are measured by selecting a region in data that
is significantly enriched in fake electrons while still kinematically close
enough to the measurement region. This is not easily achieved because fake electrons
can very closely mimic the attributes of prompt electrons, and the contribution
of prompt leptons cannot be completely suppressed, so one needs to explicitly
subtract the expected prompt electron contamination using MC predictions.


\subsection{Electron Fake Factor Measurement}

To select a region with a significant content of fake electrons, a special
single electron region is defined. The cut on \( \met + \mtW \) has been
reversed to require the observable smaller than \qty{60}{\GeV}. Additionally,
to increase the statistics, jet and \bjet selections have been loosened
to at least 3 jets of which at least one should be a \bjet. The region requirements
are summarised in \cref{tab:FakeRegionDefinitions}.

\begin{table}[htbp]
  \begin{center}
    \begin{tabular}{l c c}
      \toprule
      Selection            & Electron fakes region             & Muon fakes region  \\
      \midrule
      \( N \)(central jet) & \multicolumn{2}{c}{\( \geq 3\)}                        \\
      \( N_{\bjets}\)      & \multicolumn{2}{c}{\( \geq 1 \)}                       \\
      \( \met + \mtW \)    & \multicolumn{2}{c}{\( (0, 60) \)}                      \\
      \( N \)(lep)         & 1                                 & 1                  \\
      \( N(\ell) \)        & 1                                 & 0                  \\
      \( N(\mu) \)         & 0                                 & 1                  \\
      \( \pt(\ell) \)      & \( (30, \infty) \)                & \( (30, \infty) \)                \\
      \( \pt(j1)\)         & ---                               & \( (35, \infty) \) \\
      \( \Delta \phi(\mu, \mathrm{jet})\) & --- & \( \geq 2.7\) \\
      \bottomrule
    \end{tabular}
  \end{center}
  \caption{A summary of the event selection for fake measurement regions.}%
  \label{tab:FakeRegionDefinitions}
\end{table}

All object definitions are the same as defined in \cref{ch:definitions}.
The defined region is called the \textit{fakes-enriched region}.
Two sub-regions are defined out of the fakes-enriched region, \textit{tight}
and \textit{loose}, which are filled with electrons satisfying the corresponding
tightness requirement.

\begin{figure}[htbp]
  \centering
  \includegraphics[width=0.4\textwidth]{fakes/fake_factors_e_nominal}%
  \label{fig:Fakes_Electrons-nominal}
  \caption{Electron fake factors measured as a function of \pt for different \eta bins.
  }%
  \label{fig:Fakes_Electrons}
\end{figure}

\begin{figure}[htbp]
  \centering
  \subfloat[]{
    \includegraphics[page=7, width=0.4\textwidth]{fakes/Fakes_Electrons_e_nominal}%
    \label{fig:Fakes_Electrons-pt-loose}
  }
  \subfloat[]{
    \includegraphics[page=3, width=0.4\textwidth]{fakes/Fakes_Electrons_e_nominal}%
    \label{fig:Fakes_Electrons-pt-tight}
  }
  \\
  \subfloat[]{
    \includegraphics[page=12, width=0.4\textwidth]{fakes/Fakes_Electrons_e_nominal}%
    \label{fig:Fakes_Electrons-eta-loose}
  }
  \subfloat[]{
    \includegraphics[page=11, width=0.4\textwidth]{fakes/Fakes_Electrons_e_nominal}%
    \label{fig:Fakes_Electrons-eta-tight}
  }

  \caption{Fakes-enriched regions in the nominal selection.
    \protect\subref{fig:Fakes_Electrons-pt-loose} \pt distribution of loose electrons in the fakes-enriched region,
    \protect\subref{fig:Fakes_Electrons-pt-tight} \pt distribution of tight electrons in the fakes-enriched region,
    \protect\subref{fig:Fakes_Electrons-eta-loose} \eta distribution of loose electrons in the fakes-enriched region,
    \protect\subref{fig:Fakes_Electrons-eta-tight} \eta distribution of tight electrons in the fakes-enriched region.
    All the distributions show data events and the prompt MC component subtracted from data, to ensure a fake dominated region.
  }%
  \label{fig:Fakes_Electrons-distributions}
\end{figure}

The fakes-enriched regions are shown in \cref{fig:Fakes_Electrons-distributions}.
Selected events still contain prompt electrons after all selection requirements
are applied. These residual prompt electrons are subtracted from data before
the calculation of fake factors using the Monte Carlo simulation.
The MC samples used are \Wjets (before the fit), \Zjets, \ttbar and single top, diboson, and rare
top decays. The MC subtraction is much larger in the tight region
compared to the loose region and the number of subtracted prompt leptons
in that region amounts to up to \qty{100}{\%} of all electrons.
The \pt dependence of the fake factor is shown in \cref{fig:Fakes_Electrons}.
The fake factor was measured in three \eta slices (two slices for the forward
region and one slice for the barrel region). The last \pt bin includes the
overflow and is extended to infinity and no extrapolation of the fake
factor is performed.


\subsection{Systematic Uncertainties on Electron Fake-factors}%
\label{sec:bkg-fakes-electron-sys}

The main sources of the systematic uncertainties of the fake factor arise from
Monte Carlo modelling of the subtracted leptons, from different composition
of fake electrons in the fake enriched region compared to the measurement region,
and from the normalisation of Monte Carlo samples in the fakes-enriched region.

Given that the amount of tight fakes is very low, only the MC normalisation
uncertainty is evaluated by varying all samples simultaneously by \qty{10}{\%} to account
for cross-section and luminosity uncertainties as well as the modelling of
the prompt MC used in the subtraction procedure.

Additionally the effect of the \( \met + \mtW \) cut is tested by varying it
to \qty{50}{\GeV} and \qty{70}{\GeV}. The effect of this variation on the fake
factors is negligible.

The effect on the number of \(b-\)jets was tested by varying \( N_{\bjets}\) to include at least 2 \(b-\)jets. This effect isn't negligible and can be quite large in some bins in the electron \pt distribution. 

This region is kinematically quite far from the measurement region so the observed
variation in the fake factor due to varying composition of the fakes enriched
sample is a quite conservative estimate of the systematic uncertainty.

% \begin{table}[htbp]
%   \begin{center}
%     \begin{tabular}{ c c }
%       \toprule
%       Variation & Purpose \\
%       \midrule
%       Flipped requirement on number of jets      & fakes composition                 \\
%       Removed \met requirement      & fakes composition                 \\
%       MC scaled up by \qty{10}{\%}     & MC modelling, cross-section and luminosity \\
%       MC scaled down by \qty{10}{\%}   & MC modelling, cross-section and luminosity \\
%       \bottomrule
%     \end{tabular}
%   \end{center}

%   \caption{Summary of the variations used for the determination of the systematic uncertainty of the fake factors.}%
%   \label{tab:fakes-electron-sys}
% \end{table}

The resulting fake factors are shown on \crefrange{fig:Fakes_Electrons-elsys1}{fig:Fakes_Electrons-elsys4}.
Combined systematics of the fake factors was calculated by adding all variations
and the statistical uncertainty in quadrature. The measurement is relatively
stable overall up to very high electron \pt but the maximum uncertainty can get
very high, up to \qty{100}{\%}.

\begin{figure}[htbp]
  \centering
  \subfloat[]{
    \includegraphics[page=1, width=0.4\textwidth]{fakes/fake_factors_e_final}%
    \label{fig:Fakes_Electrons-elsys1}
  }
  \subfloat[]{
    \includegraphics[page=2, width=0.4\textwidth]{fakes/fake_factors_e_final}%
    \label{fig:Fakes_Electrons-elsys2}
  }
  \\
  \subfloat[]{
    \includegraphics[page=3, width=0.4\textwidth]{fakes/fake_factors_e_final}%
    \label{fig:Fakes_Electrons-elsys3}
  }
  \subfloat[]{
    \includegraphics[page=4, width=0.4\textwidth]{fakes/fake_factors_e_final}%
    \label{fig:Fakes_Electrons-elsys4}
  }

  \caption{The measured fake factor with systematic uncertainties.
    \protect\subref{fig:Fakes_Electrons-elsys1} fake factor variations for the first \eta bin,
    \protect\subref{fig:Fakes_Electrons-elsys2} fake factor variations for the second \eta bin,
    \protect\subref{fig:Fakes_Electrons-elsys3} fake factor variations for the third \eta bin and
    \protect\subref{fig:Fakes_Electrons-elsys4} fake factor variations for the fourth \eta bin.
  }
\end{figure}


\subsection{Validation of Fake-factors for Electrons}%
\label{subsec:bkg-fakes-electron-validation}

The measured fake factors were validated by performing the direct closure test
in the fakes-enriched region. The same event selection was used as for the
estimation, but no special binning has been applied. Fakes have been estimated
from data using the measured fake factors. The reconstructed object selection
was the same as the nominal event selection (\cref{ch:definitions}).

Prompt Monte Carlo processes that were used are the same as for the estimation.
The remaining events are described well by the fake factor method as seen on
\cref{fig:Fakes_ClosureElectrons}.
The only considered systematic uncertainty in this region was
the systematic uncertainty of the fake factor. Overall the agreement is good
even up to very high electron energies.

\begin{figure}[htbp]
  \centering
  \subfloat[]{
    \includegraphics[page=2, width=0.4\textwidth]{fakes/Fakes_ClosureElectrons_e_validation_sys}%
    \label{fig:Fakes_ClosureElectrons-pt}
  }
  \subfloat[]{
    \includegraphics[page=5, width=0.4\textwidth]{fakes/Fakes_ClosureElectrons_e_validation_sys}%
    \label{fig:Fakes_ClosureElectrons-eta}
  }
  \\
  \subfloat[]{
    \includegraphics[page=7, width=0.4\textwidth]{fakes/Fakes_ClosureElectrons_e_validation_sys}%
    \label{fig:Fakes_ClosureElectrons-met}
  }
  \subfloat[]{
    \includegraphics[page=13, width=0.4\textwidth]{fakes/Fakes_ClosureElectrons_e_validation_sys}%
    \label{fig:Fakes_ClosureElectrons-mt}
  }

  \caption{The electron fake-factors closure.
    \protect\subref{fig:Fakes_ClosureElectrons-pt} electron \pt distribution,
    \protect\subref{fig:Fakes_ClosureElectrons-eta} electron \eta distribution,
    \protect\subref{fig:Fakes_ClosureElectrons-met} \met distribution, and
    \protect\subref{fig:Fakes_ClosureElectrons-mt} \mtW distribution.
    The dashed band shows the systematic uncertainty associated to the fake factor measurement.
  }%
  \label{fig:Fakes_ClosureElectrons}
\end{figure}


\pagebreak

% Fake Background Muons
\section{Muon Fake Factor Measurement}%
\label{sec:bkg-fakes-muon}

Fake muons arise from in-flight decays of mesons inside jets with
the semi-leptonic decay of B meson as one of the largest contributions.
Combinations of prompt and fake muons can contribute to the measurement region
if the fake muons satisfy tight isolation and impact parameter requirements.
The muon fake contribution to the measurement region is estimated as described
in \cref{sec:bkg-fakefactor}, which involves a data driven estimate
of the fake factor.


\subsection{Muon Fake Factor Measurement}

Fake-factors are measured requiring a jet to be associated to a muon, thereby
targeting a fake factor measurement in di-jet events. Basing on the selection defined
for electrons in \cref{tab:FakeRegionDefinitions}, events are selected
by requiring the highest \pt light jet in the event to have an azimuthal angular distance
greater than \( \Delta \phi > 2.7 \) to the muon. This topology is qualitatively illustrated
in \cref{fig:bkg-dijet}.

\begin{figure}[htbp]
  \begin{center}
    \includegraphics[width=0.6\textwidth]{fakes/dijet}
    \caption{Qualitative sketch of a di-jet event selected for the measurement
      of the muon fake-factors.}%
    \label{fig:bkg-dijet}
  \end{center}
\end{figure}

The selection of di-jet events allows one to achieve high statistics for
measuring fake-factors, while also allowing to address systematic uncertainties
using the properties of the recoiling tag jet.

The muon and the leading=\pt jet in the event have to be back-to-back in the transverse
plane, i.e. \( \Delta \phi(\mu, \mathrm{jet}) > 2.7 \) and have \pt threshold
of \( \pt(\mu) > \qty{30}{\GeV} \) and \( \pt(\mathrm{jet}) > \qty{35}{\GeV} \).
The selection criteria of the muon fake measurement region are summarised in \cref{tab:FakeRegionDefinitions}.

The distributions of muon \pt and \eta after all selection requirements
are applied, are reported in \cref{fig:Fakes_Muons-distributions}, where
tight and loose muons are tight and loose, respectively, as defined in
\cref{tab:definitions-muons}. The \pt dependence on fake factors is shown
in \cref{fig:Fakes_Muons}.

\begin{figure}[htbp]
  \centering
  \includegraphics[width=0.4\textwidth]{fakes/fake_factors_m_nominal}
  \caption{Muon fake factors measured as a function of \pt.}%
  \label{fig:Fakes_Muons}
\end{figure}

\begin{figure}[htbp]
  \centering
  \subfloat[]{
    \includegraphics[page=3, width=0.4\textwidth]{fakes/Fakes_Muons_m_nominal}%
    \label{fig:Fakes_Muons-pt-loose}
  }
  \subfloat[]{
    \includegraphics[page=1, width=0.4\textwidth]{fakes/Fakes_Muons_m_nominal}%
    \label{fig:Fakes_Muons-pt-tight}
  }
  \\
  \subfloat[]{
    \includegraphics[page=6, width=0.4\textwidth]{fakes/Fakes_Muons_m_nominal}%
    \label{fig:Fakes_Muons-eta-loose}
  }
  \subfloat[]{
    \includegraphics[page=5, width=0.4\textwidth]{fakes/Fakes_Muons_m_nominal}%
    \label{fig:Fakes_Muons-eta-tight}
  }

  \caption{Fakes-enriched regions in the nominal selection.
    \protect\subref{fig:Fakes_Muons-pt-loose} \pt distribution of loose muons in the fakes-enriched region,
    \protect\subref{fig:Fakes_Muons-pt-tight} \pt distribution of tight muons in the fakes-enriched region,
    \protect\subref{fig:Fakes_Muons-eta-loose} \eta distribution of loose muons in the fakes-enriched region and
    \protect\subref{fig:Fakes_Muons-eta-tight} \eta distribution of tight muons in the fakes-enriched region.
    All the distributions show data events and the prompt MC component subtracted from data, to ensure a fake dominated region.
  }%
  \label{fig:Fakes_Muons-distributions}
\end{figure}


\subsection{Systematic Uncertainties on Muon Fake Factors}%
\label{sec:bkg-fakes-muon-sys}

Similar to electrons all Monte Carlo samples were varied by \qty{10}{\%}
to account for cross-section and luminosity uncertainties as well as the
modelling of the prompt MC used in the subtraction procedure.

Additionally the effect of the \( \met + \mtW \) cut is tested by varying it
to \qty{50}{\GeV} and \qty{70}{\GeV}. The effect of this variation on the fake
factors is not negligible.

The effect on the number of \bjets was tested by varying \( N_{\bjets}\) to
include at least 2 \bjets. This effect is not negligible and can be quite large
in some bins in the electron \pt distribution. 

Additional variations were tested, including the variation of the transverse
momentum of the leading jet, with thresholds of \( \pt(j1) > \qty{40}{\GeV} \)
and \qty{30}{\GeV}, \( \Delta \phi(\mu, \mathrm{jet})\) > 2.6 and 2.8,
and the effect of omitting the requirement on light jets.

The effect of each systematic alteration on the nominal measurement can be
seen on \crefrange{fig:Fakes_Muons-musys1}{fig:Fakes_Muons-musys3}.
The total uncertainty is estimated by comparing the statistical uncertainty on
the nominal measurement with the maximum deviation between the nominal
fake-factor and each systematically altered measurement. The largest signed
deviation is taken to be the total uncertainty for lower bounds and upper bounds
in turn.

\begin{figure}[htbp]
  \centering
  \subfloat[]{
    \includegraphics[page=1, width=0.4\textwidth]{fakes/fake_factors_m_final}%
    \label{fig:Fakes_Muons-musys1}
  }
  \subfloat[]{
    \includegraphics[page=2, width=0.4\textwidth]{fakes/fake_factors_m_final}%
    \label{fig:Fakes_Muons-musys2}
  }
  \\
  \subfloat[]{
    \includegraphics[page=3, width=0.4\textwidth]{fakes/fake_factors_m_final}%
    \label{fig:Fakes_Muons-musys3}
  }

  \caption{The measured fake factor for each systematics variation.
    \protect\subref{fig:Fakes_Muons-musys1} fake factor variations for the first \eta bin,
    \protect\subref{fig:Fakes_Muons-musys2} fake factor variations for the second \eta bin and
    \protect\subref{fig:Fakes_Muons-musys3} fake factor variations for the third \eta bin.
  }% 
  \label{fig:Fakes_Muons-musys}
\end{figure}


\subsection{Validation of Fake-factors for Muons}%
\label{subsec:bkg-fakes-muon-validation}

The measured fake factors were validated by performing the direct closure test
in the fakes-enriched region. The same event selection was used as for the
estimation, but no special binning has been applied. Fakes have been estimated
from data using the measured fake factors. The reconstructed object selection
was the same as the nominal event selection (\cref{ch:definitions}).

Prompt Monte Carlo processes that were used are \Wjets, \Zjets, \ttbar and
single top, diboson, and rare top decays. The remaining events
are described well by the fake factor method as seen on
\cref{fig:Fakes_ClosureMuons}.
The only considered uncertainty in this region is the systematic uncertainty
of the fake factor.

\begin{figure}[htbp]
  \centering
  \subfloat[]{
    \includegraphics[page=3, width=0.4\textwidth]{fakes/Fakes_ClosureMuons_m_validation_sys}%
    \label{fig:Fakes_ClosureMuons-pt}
  }
  \subfloat[]{
    \includegraphics[page=7, width=0.4\textwidth]{fakes/Fakes_ClosureMuons_m_validation_sys}%
    \label{fig:Fakes_ClosureMuons-eta}
  }
  \\
  \subfloat[]{
    \includegraphics[page=9, width=0.4\textwidth]{fakes/Fakes_ClosureMuons_m_validation_sys}%
    \label{fig:Fakes_ClosureMuons-met}
  }
  \subfloat[]{
    \includegraphics[page=15, width=0.4\textwidth]{fakes/Fakes_ClosureMuons_m_validation_sys}%
    \label{fig:Fakes_ClosureMuons-mt}
  }

  \caption{The muon fake-factors closure.
    \protect\subref{fig:Fakes_ClosureMuons-pt} muon \pt distribution and
    \protect\subref{fig:Fakes_ClosureMuons-eta} muon \eta distribution,
    \protect\subref{fig:Fakes_ClosureMuons-met} \met distribution, and
    \protect\subref{fig:Fakes_ClosureMuons-mt} \mtW distribution.
    The dashed band shows the systematic uncertainty associated to the fake factor measurement.
  }%
  \label{fig:Fakes_ClosureMuons}
\end{figure}

